\subsection{停止、观察资格与路径退出}
\label{sec:process-theory}

\subsubsection{从重建片段到实际观察路径}
完整片段按第~\ref{subsec:segment-rule}~节的高置信候选修改边连接记录；实际观察路径则是按第~\ref{subsec:failure-cohort}~节停止规则截取的部分。一步是两个送审版本之间的转移，而不是一次内部编辑、模型调用或一段固定时间。路径由保存的记录重建，同一任务可以有几条彼此分开的路径。

假设记录依次为失败、失败、通过。从第一个失败起点出发，共有两步，观察在通过时停止。因此，三份候选记录给出了两次转移。链条也可能在最新审核尚未通过时结束，这称为\emph{路径退出}。任务以后仍可能在别处继续，退出因而不表示永久失败。第~\ref{sec:process}~节将按这个观察规则研究历史数据；这里先看它在简单概率模型中的后果。

模型中的成功指记录通过。\emph{未通过}包括其余所有标签，也包括部分完成或未决判断；失败起点路径则必须从明确的失败标签开始。这些约定描述审核记录，而不是独立认证的题意符合性。

\subsubsection{恒定条件通过概率与可预测观察}
考虑反复投掷一枚正面概率为 $p$ 的硬币，其中 $0<p<1$，直到首次正面出现。如果用正面表示记录通过，就得到一个每步通过概率恒定的连续送审模型。真实候选可能改善，审核条件也可能变化；这个说明性模型暂时把这些可能性放在一边。

为准确写出计算需要的条件，令 $j=1,2,\ldots$ 表示一个可能的步骤。用 $H_{j-1}$ 表示\emph{得知第 $j$ 步结果之前}已有的全部信息，例如先前候选、先前审核结果，以及已经作出的是否继续观察的决定。定义观察指示变量、接受指示变量及其累计计数：
\begin{center}
\begin{tabularx}{\linewidth}{lY}
\toprule
符号 & 对一条路径而言的含义 \\
\midrule
$B_j$ & 如果此前没有通过，而且第 $j$ 步将被观察，则为 1；否则为 0。 \\
$A_j$ & 如果这一步的审核记录通过，则为 1；未通过则为 0。当 $B_j=0$ 时也把它记为 0，只是为了让记号完整。 \\
$N_k=\sum_{j=1}^k B_j$ & 截至第 $k$ 步实际观察到的合格步骤数。 \\
$D_k=\sum_{j=1}^k B_jA_j$ & 这些步骤中记录通过的次数。由于第一次通过就停止，它至多为 1。 \\
\bottomrule
\end{tabularx}
\end{center}

说 $B_j$ 可以根据 $H_{j-1}$ 确定，是指在看见这一步的结果之前，就能决定这条观察是否进入分析。这种性质常称为\emph{可预测的观察规则}。它不是说我们能够预测下一份候选会通过。例如，“未通过就继续，已经通过就停止”只使用已有信息；“只有下一份审核看起来有希望时才保留这一行”则可能使用尚未出现的信息，不一定满足这个条件。

对于合格步骤，恒定概率假设写成
\begin{equation}
 P(A_j=1\mid H_{j-1})=p\qquad\text{只要 }B_j=1.
 \label{eq:constant-conditional-pass}
\end{equation}
这要求在每一种合格过去历史之后具有相同的条件通过概率，比汇总通过比例相等更强。等价地，在合格历史上有 $E(A_j\mid H_{j-1})=p$。

\begin{proposition}[明确观察规则下的恒定条件概率]
\label{prop:stopping}
假设在得知第 $j$ 步结果之前就已知道观察资格 $B_j$，并且式~\eqref{eq:constant-conditional-pass}~对每一种合格历史都成立。那么，对于每个预先固定的有限步数 $k$，
\begin{equation}
 E(D_k-pN_k)=0,\qquad\text{从而}\qquad E(D_k)=pE(N_k).
 \label{eq:stopped-count-identity}
\end{equation}
进一步，如果第 1 步会被观察，并且每次未通过之后都继续观察下一步，直到第一次通过，那么首次通过的步数 $T$ 满足
\begin{equation}
 P(T=j)=(1-p)^{j-1}p,\qquad j=1,2,\ldots.
 \label{eq:geometric-pass}
\end{equation}
\end{proposition}

\begin{proof}
在合格步骤中，已知此前历史时，$A_j-p$ 的条件平均值为 $p-p=0$。在不合格步骤中，$B_j(A_j-p)=0$。由于观察资格已知，这两种情形都给出
\[
 E\{B_j(A_j-p)\mid H_{j-1}\}=0.
\]
再对可能的历史取平均，结果仍为零。对固定的步骤 $1,\ldots,k$ 把这些等式相加，期望内的总和恰好是 $D_k-pN_k$，于是得到式~\eqref{eq:stopped-count-identity}。

如果每次未通过之后都继续观察，那么第一次在第 $j$ 步通过，需要先连续 $j-1$ 次未通过，然后通过。把每一步的条件概率相乘，得到
\[
 \underbrace{(1-p)\cdots(1-p)}_{j-1\text{ 次未通过}}p.
\]
这就是式~\eqref{eq:geometric-pass}。较早的步数通常获得更大权重，因为要到更晚才首次通过，必须先经历此前的全部未通过。这个计算使用已经声明的条件概率；在这个涵盖全部过去信息的条件之外，不需要再额外添加一个独立性假设。
\end{proof}

式~\eqref{eq:geometric-pass}~给出几何分布的首次通过步数 $T$，它与任务符合性参照 $T(x)$ 不同。例如 $p=0.2$ 时，第 1、2、3 步首次通过的概率分别为 $0.2$、$0.8(0.2)=0.16$、$0.8^2(0.2)=0.128$。尽管每一步剩下的合格路径越来越少，在这个例子中，\emph{仍然合格的路径当中}下一步通过的概率始终为 $0.2$。因此，合格路径集合缩小，本身不能否定恒定概率模型；当然也不能据此认定历史数据符合这个模型。

\subsubsection{鞅表述}
\label{subsec:martingale}
早期分析有时使用“鞅”这个术语。在这里，它描述的是累计偏差
\[
 M_k=D_k-pN_k=\sum_{j=1}^k B_j(A_j-p).
\]
已知此前信息之后，它下一步变化的平均值为零，即 $E(M_k\mid H_{k-1})=M_{k-1}$。这就是这里真正需要的数学性质，上面已经证明。“过滤族”只是不断增加的信息集合的技术名称，这里用历史信息 $H_0,H_1,\ldots$ 来表达。因此，鞅性质是对上面零均值增量计算的另一种表述。

\subsubsection{随机分母下的平均}
式~\eqref{eq:stopped-count-identity}~讨论的是\emph{计数}的期望。它没有说 $E(D_k/N_k)=p$，因为实际观察的步数本身可以依赖此前的结果。一个两步硬币例子就能看清区别：取 $p=1/2$，第一次通过就停止，最多观察两步。
\begin{center}
\begin{tabular}{lrrrr}
\toprule
观察到的序列 & 概率 & $D_2$ & $N_2$ & $D_2/N_2$ \\
\midrule
通过 & $1/2$ & 1 & 1 & 1 \\
未通过、通过 & $1/4$ & 1 & 2 & $1/2$ \\
未通过、未通过 & $1/4$ & 0 & 2 & 0 \\
\bottomrule
\end{tabular}
\end{center}
这里 $E(D_2)=3/4$、$E(N_2)=3/2$，所以 $E(D_2)/E(N_2)=1/2$。但是，各条最终路径比例的平均值为 $E(D_2/N_2)=1/2+1/8=5/8$。快速成功的路径只经过一次观察，就贡献一个等于 1 的比例。因此，交换“先求比例”和“先取期望”的顺序会改变答案。

如果不设步数上限、一直观察下去，用 $D$ 和 $N$ 表示最终停止时的\emph{全部}通过次数和已观察步数；它们不同于只累计到某个固定有限步骤 $k$ 的 $D_k,N_k$。由于 $P(T>k)=(1-p)^k$ 趋于零，第一次通过会以概率 1 出现。因此，以概率 1 有 $D=1$、$N=T$，那个概率为零的例外不影响下面的期望。同样的区别可以写成一个精确公式：
\[
 E(D/N)=\sum_{j=1}^{\infty}\frac{p(1-p)^{j-1}}{j}
 =\frac{-p\log p}{1-p}>p.
\]
把几何级数 $\sum_{r=0}^{\infty}x^r=1/(1-x)$ 从 0 积分到 $1-p$，得到 $\sum_{j\ge1}(1-p)^j/j=-\log p$。又因为 $0<p<1$ 时有 $-\log p=\int_p^1(1/u)\,du>1-p$，严格不等式随之成立。这是说明性模型的性质，不是本报告用于校正历史数据的公式。

\subsubsection{退出与成功路径筛选}
没有退出的硬币模型假设每次未通过之后都还有下一次观察。现在改成：每次未通过之后，再用第二枚独立硬币决定路径是否继续。令 $a$ 表示继续的概率，$0\le a\le1$。在这个说明性构造中，所有决定继续的硬币和决定通过的硬币彼此独立。一个合格步骤之后，就有下列三种结果：
\begin{center}
\begin{tabularx}{\linewidth}{Yl}
\toprule
发生的事情 & 概率 \\
\midrule
记录通过，然后停止 & $p$ \\
记录未通过，然后退出路径 & $(1-p)(1-a)$ \\
记录未通过，然后继续观察下一步 & $(1-p)a$ \\
\bottomrule
\end{tabularx}
\end{center}
三个概率相加等于 1。把第三个概率简写为 $b=(1-p)a$：它表示既在当前步骤未通过，\emph{又}能进入下一次观察的概率。它不是通过概率。

要在第 $j$ 步观察到通过，前 $j-1$ 步必须每次都出现第三种结果，然后第 $j$ 步通过。因此，
\[
 P(\text{退出之前在第 }j\text{ 步通过})=b^{j-1}p.
\]
不同 $j$ 对应的首次通过事件不会同时发生。把它们的概率相加，再用 $1+b+b^2+\cdots=1/(1-b)$，得到
\begin{equation}
 s:=P(\text{退出之前通过})=\frac{p}{1-b}.
 \label{eq:pass-before-exit}
\end{equation}
这里 $s$ 是关于\emph{整条路径}的概率；$p$ 是关于一个合格步骤的概率。它们的分母不同，回答的也是不同问题。

最后，假设我们删去所有未通过就退出的路径，只报告成功路径。令 $J$ 表示这些路径记录通过的步数，使用普通的条件概率公式便有
\begin{equation}
 P(J=j\mid\text{退出之前通过})
 =\frac{p b^{j-1}}{s}=(1-b)b^{j-1}.
 \label{eq:selected-success-time}
\end{equation}
这又是一个几何分布，但它的参数是 $1-b=p+(1-p)(1-a)$；当 $a<1$ 时，这个数大于 $p$。只保留成功路径的筛选偏向较短历史：要成为一条较长的成功路径，记录必须先避开更多次退出机会。因此，看起来更快的首次通过分布，并不证明每一步通过的概率更高。

下面给出一个具体的辨识问题。考虑两种可能的机制：
\begin{center}
\begin{tabular}{lrrrr}
\toprule
机制 & $p$ & $a$ & $b=(1-p)a$ & $s=p/(1-b)$ \\
\midrule
甲 & 0.2 & 0.5 & 0.4 & $1/3$ \\
乙 & 0.5 & 0.8 & 0.4 & $5/6$ \\
\bottomrule
\end{tabular}
\end{center}
两者在成功路径中的通过步数概率完全一样，依次为 $0.6$、$0.24$、$0.096$ 等等。但是，它们每个合格步骤的通过概率分别为 $0.2$ 和 $0.5$。只看成功路径的长度，就无法分清是哪一种机制。总体通过比例则能把它们区分开：如果同时知道 $s$ 和 $b$，式~\eqref{eq:pass-before-exit}~会给出 $p=s(1-b)$。这里的不确定性来自保留了哪些信息，不是概率计算本身失效。

现在可以区分三个量：合格步骤的通过率、固定路径的终态，以及筛选出的成功路径长度。若要把模型用于历史记录，还需核查它的观察假设。事后划分修改边可能用到下一候选或审核的信息，因此复现这种划分，并不能证明命题~\ref{prop:stopping}~要求的结果出现之前决定观察资格这一条件。后面的实证分析因而限于保留路径及其观察关联。